\section{Conclusion}

This study presented a hybrid retrieval and reranking framework for citation-aware retrieval-augmented generation. The proposed pipeline used Amazon Bedrock Knowledge Base for document ingestion, parsing, chunking, embedding generation, and evidence retrieval. Source PDF documents were stored in Amazon S3, embedded using Amazon Titan Text Embeddings V2, and indexed using Amazon OpenSearch Serverless. A hybrid retrieval strategy was then combined with reranking to prioritize the most relevant evidence chunks before answer generation.

The experimental results demonstrate that the proposed framework produced highly grounded responses for the evaluated biomedical natural language processing and healthcare transformer queries. Across 25 evaluation queries, the system retrieved and reranked 500 evidence chunks and generated responses using the top-ranked evidence. Claim-level grounding evaluation extracted 200 factual claims from the generated answers. All extracted claims were judged to be supported by the retrieved evidence, resulting in an overall grounding accuracy of 100.0\%. These findings indicate that the combination of hybrid retrieval, reranking, conservative prompting, and independent claim-level evaluation can improve the reliability of RAG-generated responses when sufficient evidence is available in the source document collection.

The reranking stage also played an important role in evidence prioritization. Only a small portion of the retrieved chunks retained the same rank after reranking, indicating that the reranker substantially changed the initial retrieval order. This suggests that reranking can help identify more directly relevant evidence from the candidate chunks returned by the initial retrieval stage. The distribution of retrieved evidence across multiple source documents further shows that the framework did not depend on a single document but was able to draw support from different papers within the knowledge base.

A key contribution of this work is the use of claim-level grounding evaluation rather than only answer-level assessment. By decomposing each generated response into factual claims, the evaluation process provided a more detailed measure of whether individual statements were supported by retrieved evidence. The use of a separate judge model for grounding assessment also helped reduce potential self-evaluation bias by separating the answer-generation model from the evaluation model.

However, the current study has several limitations. Because the present work was designed as a pilot-scale proof-of-concept study, the results should be interpreted as an initial feasibility evaluation rather than as a large-scale benchmark comparison. First, the evaluation was conducted on a relatively small set of 25 queries, which limits the generalizability of the findings. Second, the reported accuracy represents internal grounding accuracy rather than strict citation accuracy, because the generated answers did not include explicit source labels for each factual claim. Third, the grounding judgments were produced by an automated judge model and were not independently verified by human domain experts. Finally, the experiment did not include a direct comparison with retrieval-only, reranking-only, or non-hybrid baseline systems.

Future work should expand the evaluation to a larger and more diverse query set across multiple domains and document collections. Additional experiments should compare hybrid retrieval with purely lexical, purely vector-based, and retrieval-only baselines. The framework should also be extended to require explicit sentence-level source citations in the generated answers, allowing strict citation accuracy to be evaluated in addition to general grounding accuracy. Human expert review should be incorporated to validate automated grounding judgments and assess whether the generated responses are not only evidence-supported but also clinically and scientifically appropriate.

Overall, the results suggest that a hybrid retrieval and reranking approach can provide a strong foundation for reliable and evidence-grounded RAG systems. By combining structured document ingestion, semantic retrieval, reranking, controlled answer generation, and claim-level grounding evaluation, the proposed framework offers a reproducible approach for improving the trustworthiness of generated responses in biomedical and healthcare-related information retrieval tasks.

\section*{Reproducibility and Data Availability}

The materials required to reproduce the experimental workflow, including the query set, retrieved and reranked evidence outputs, generated answers, claim-level grounding evaluation results, query-level accuracy summaries, and overall summary files, are available in the project repository at: \url{https://drive.google.com/drive/folders/1QLNxniZ_Zxp42p8Y5fC_XQ6sWJzio9lP?usp=sharing}. The implementation was developed in Python using Google Colab and external managed services, including Amazon Bedrock Knowledge Bases, Amazon OpenSearch Serverless, Amazon Bedrock Runtime, and Cohere reranking. To support reproducibility, intermediate outputs were stored as structured CSV files so that retrieval, reranking, answer generation, and grounding evaluation could be inspected independently. API credentials, access keys, and service-specific authentication information are not included in the shared files for security reasons and should be supplied through secure environment variables when reproducing the experiments.

\section*{AI Use Disclosure}
The authors used large language models as experimental components within the proposed RAG pipeline. No AI tool was used to fabricate data or alter experimental results. The authors reviewed and verified the manuscript content, tables, and reported findings.